\HMTNarrativeHeading{Proportion, raffinement et restitution}

L’extrême et moyenne raison holographique réunit ici la complétion
arithmétique et la récupération opératorielle. Le raffinement avec retenue
redistribue une unité entre composantes et conserve la somme
normalisée. La division euclidienne, appliquée avec la profondeur et
la convention de retenue déclarées, récupère les étiquettes
antérieures. La conservation et l’inverse reçoivent ainsi des opérations
concrètes.

La norme ternaire apporte l'entier soixante-treize et son interprétation algébrique et réticulaire. La relation avec la complétion décimale s'exprime au moyen d'identités exactes. Le bilan bilatéral étend l'analyse à deux transports isométriques entre espaces de dimension arbitraire. Les termes croisés s'annulent dans une combinaison pondérée et restituent le carré de la norme initiale.

La spécialisation nonadique détermine la répartition initiale des poids
soixante-douze sur soixante-treize et un sur soixante-treize.
Un mélange orthogonal produit ensuite les deux canaux normalisés.
Lorsque les transports admettent le relatif unitaire déclaré,
le lecteur effectif combine l’observable avec son conjugué ; son
invariance équivaut alors à la commutation avec ce relatif. Le transport
de l’état complet conserve simultanément les relations
que chaque observable permet d’évaluer.

La politique d’archivage prolonge ce bilan. Les mémoires
successives et le terminal maintiennent la totalité à profondeur
finie. L’estimation uniforme du terminal permet d’obtenir
une archive infinie récupérable au moyen de l’adjoint de l’opérateur
isométrique. Les bornes décrivent la stabilité de cette récupération
lorsque la profondeur augmente.

La constante de couplage participe à la même construction. Son orbite et sa forme de Gram permettent d'identifier des opérateurs ; l'isométrie d'archivage conserve ces réponses et leurs bornes. Le raffinement des étiquettes est relevé en un transport unitaire et se compose avec le bilan bilatéral. L'unité conserve ainsi son contenu arithmétique, sa réalisation linéaire et la possibilité de restitution. Tel est l'aboutissement opératoire de l'argument directeur : une structure engendrée admet des transformations qui redistribuent sa lecture et conservent les données de sa récupération.
